% p015_c 的电脑重画（旋转圆：入射点 O 处 V0 向上，两个轨迹圆①②，圆心圆，远点圆；表格文字不在图内）
\begin{tikzpicture}[line width=0.7pt, >={Stealth[length=2.2mm]}]
  % 远点圆（以 O 为心、半径 2r，被图的上下边界截断）与左侧相切的淡圆
  \begin{scope}
    \clip (-6.8,-2.7) rectangle (6.4,2.2);
    \draw[gray!70, line width=0.5pt] (0,0) circle (3.4);
  \end{scope}
  \draw[gray!70, line width=0.5pt] (-4.86,-0.61) circle (1.5);
  % 圆心圆（圆心的轨迹，以 O 为心、半径 r）
  \draw[gray!70, line width=0.5pt] (0,0) circle (1.7);
  % 两个轨迹圆
  \draw (-1.7,0) circle (1.7);
  \draw (1.7,0) circle (1.7);
  % 三个点：圆心、入射点、圆心
  \fill (-1.7,0) circle (1.8pt) (0,0) circle (1.8pt) (1.7,0) circle (1.8pt);
  % 速度
  \draw[->, thick] (0,0) -- (0,0.95);
  \node[left] at (0,0.3) {$V_0$};
  % 标注
  \node[draw, circle, inner sep=0.8pt] at (-0.4,1.95) {\footnotesize 2};
  \node[draw, circle, inner sep=0.8pt] at (1.75,1.95) {\footnotesize 1};
  \node at (1.55,0.95) {圆心圆};
  \node at (3.75,1.68) {远点圆};
  \node at (1.1,-0.8) {$r=\dfrac{mv}{qB}$};
  \node at (4.8,0.76) {$d=2r=\dfrac{2mv}{qB}$};
\end{tikzpicture}
